\section{Results}
Table~\ref{tab:results} summarizes performance across all five terrain classes.
\textsc{ContactRefine} reduces stumble rate relative to the elevation-only baseline on every
terrain class, with the largest gain on rubble ($21.4\% \to 7.1\%$), where map noise near
fragment edges is most severe. The learned-encoder baseline improves over elevation-only but
trails \textsc{ContactRefine} on all classes except wet tile, where contact-force history is
less informative since slip there is governed by surface friction rather than geometry.

\begin{table}[t]
\centering
\caption{Stumble rate (\%) and cost of transport by terrain class, averaged over $280$ trials per cell.}
\label{tab:results}
\footnotesize
\begin{tabular}{@{}lcccc@{}}
\toprule
\textbf{Terrain} & \textbf{Elev.-only} & \textbf{React.} & \textbf{Learn.-enc.} & \textbf{Ours} \\
\midrule
Rubble        & 21.4 & 15.8 & 12.9 & \textbf{7.1} \\
Loose gravel  & 18.6 & 14.2 & 10.5 & \textbf{6.6} \\
Wet tile      & 12.1 & 9.7  & \textbf{7.4}  & 7.8 \\
Staircase     & 16.3 & 12.0 & 9.8  & \textbf{4.9} \\
Root trail    & 21.0 & 16.9 & 13.7 & \textbf{4.6} \\
\midrule
Mean          & 17.9 & 13.7 & 10.9 & \textbf{6.2} \\
\bottomrule
\end{tabular}
\end{table}

Figure~\ref{fig:cot} plots mean cost of transport against stumble rate for all four methods,
pooled across terrain classes. \textsc{ContactRefine} occupies the lower-left region of the
plot, indicating it does not trade efficiency for stability -- a pattern we did not assume in
advance, since foothold corrections that avoid slip can in principle require less economical
leg trajectories.

\begin{figure}[t]
\centering
\begin{tikzpicture}
\begin{axis}[
  width=0.95\linewidth,
  height=5.1cm,
  xlabel={Stumble rate (\%)},
  ylabel={Cost of transport},
  xmin=0, xmax=20,
  ymin=0.28, ymax=0.42,
  axis lines=left,
  tick label style={font=\scriptsize},
  label style={font=\scriptsize},
  legend style={font=\scriptsize, at={(0.97,0.05)}, anchor=south east, draw=none, fill=none},
  grid=major,
  grid style={gray!15}
]
\addplot[only marks, mark=*, mark size=2.6pt, color=irosamber] coordinates {(17.9,0.401)};
\addlegendentry{Elevation-only}
\addplot[only marks, mark=square*, mark size=2.4pt, color=irosgray] coordinates {(13.7,0.379)};
\addlegendentry{Reactive}
\addplot[only marks, mark=triangle*, mark size=2.8pt, color=irosteal] coordinates {(10.9,0.352)};
\addlegendentry{Learned-encoder}
\addplot[only marks, mark=diamond*, mark size=2.8pt, color=irosnavy] coordinates {(6.2,0.336)};
\addlegendentry{Ours}
\end{axis}
\end{tikzpicture}
\caption{Cost of transport versus stumble rate, pooled across all five terrain classes.
Lower-left is better on both axes.}
\label{fig:cot}
\end{figure}

Added latency for \textsc{ContactRefine} was $1.8\,\text{ms}\pm0.3\,\text{ms}$ per swing phase
on the onboard compute, within the timing budget of the $500\,\text{Hz}$ control loop; the
learned-encoder baseline added $2.6\,\text{ms}$ owing to its larger terrain backbone.
